% SEGMENT OLP-0701-S01
% Part: proof-theory
% Chapter: sequent-calculus
% Section: invertibility

\documentclass[../../../include/open-logic-section]{subfiles}

\begin{document}

\olfileid{pt}{seq}{inv}

\olsection{விதிகளின் மீளாக்கத்தன்மை}

\begin{defn}
ஒரு முறைமையில் விதி~$R$,
\[
\AxiomC{$S_1$}
\AxiomC{$\dots$}
\AxiomC{$S_n$}
\RightLabel{$R$}
\TrinaryInfC{$S$}
\DisplayProof
\]
\emph{மீளாக்கத்தக்கது} எனப்படுவது:
$\Proves S$ என்றால், $i = 1$, \dots,~$n$
ஒவ்வொன்றுக்கும் ஏதோ ஓர் உயரத்தில்
$\Proves[h_i] S_i$ ஆக வேண்டும்.
மேலும் $\Proves_n S$ என்றால்,
$i = 1$, \dots,~$n$ ஒவ்வொன்றுக்கும்
$\Proves[n] S_i$ ஆகுமானால்,
அது அந்த முறைமையில்
\emph{!!{height}-பேணும் மீளாக்கத்தக்க}
(அல்லது \emph{!!{hp}-மீளாக்கத்தக்க})
விதி எனப்படும்.
\end{defn}

% SEGMENT OLP-0701-S02
\begin{lem}\ollabel{lem:G3c-invert}
\Log{G3c} இல் $\lnot$, $\land$, $\lor$,
$\lif$ ஆகியவற்றுக்கான விதிகள்
!!{height}-பேணும் மீளாக்கத்தக்கவை;
அதாவது:
\begin{enumerate}
  \item \RightR{\lnot}: $ \Log{G3c} \Proves[n] \Gamma \Sequent
  \Delta, \lnot !A$ எனில், $\Log{G3c} \Proves[n] !A, \Gamma \Sequent
  \Delta$.
  \item \LeftR{\lnot}: $\Log{G3c} \Proves[n] \Gamma, \lnot !A
  \Sequent \Delta$ எனில், $\Log{G3c} \Proves[n] \Gamma \Sequent \Delta,
  !A$.
  \item \RightR{\land}: $\Log{G3c} \Proves[n] \Gamma \Sequent
  \Delta, !A \land !B$ எனில், $\Log{G3c} \Proves[n] \Gamma \Sequent
  \Delta, !A$ மற்றும் $\Log{G3c} \Proves[n] \Gamma \Sequent
  \Delta, !B$.
  \item \LeftR{\land}: $\Log{G3c} \Proves[n] \Gamma, !A \land !B
  \Sequent \Delta$ எனில், $\Log{G3c} \Proves[n] !A, !B, \Gamma \Sequent
  \Delta$.
  \item \RightR{\lor}: $\Log{G3c} \Proves[n] \Gamma \Sequent
  \Delta, !A \lor !B$ எனில், $\Log{G3c} \Proves[n] \Gamma \Sequent
  \Delta, !A, !B$.
  \item \LeftR{\lor}: $\Log{G3c} \Proves[n] !A \lor !B, \Gamma \Sequent
  \Delta$ எனில், $\Log{G3c} \Proves[n] !A, \Gamma \Sequent
  \Delta$ மற்றும் $\Log{G3c} \Proves[n] !B, \Gamma \Sequent
  \Delta$.
  \item \RightR{\lif}: $\Log{G3c} \Proves[n] \Gamma \Sequent
  \Delta, !A \lif !B$ எனில், $\Log{G3c} \Proves[n] !A, \Gamma \Sequent
  \Delta, !B$.
  \item \LeftR{\lif}: $\Log{G3c} \Proves[n] !A \lif !B, \Gamma
  \Sequent \Delta$ எனில், $\Log{G3c} \Proves[n] \Gamma \Sequent \Delta,
  !A$ மற்றும் $\Log{G3c} \Proves[n] !B, \Gamma \Sequent \Delta$.
\end{enumerate}
\end{lem}

% SEGMENT OLP-0701-S03
\begin{proof}
\Log{G3c} இன் எந்த விதி~$R$ க்கும்,
அந்த $R$ இன் முடிவு~$S$ க்கு
!!a{proof}~$\pi$ இருந்தால், அதன்
அந்த $R$ இன் முன்கூற்றுகள்~$S_i$ க்கும்
!!{proof}s~$\pi_i$ இருந்து,
$\pheight{\pi_i} \le \pheight{\pi}$
ஆகும் என்பதைக் காட்ட வேண்டும்.
\LeftR{\land} விதியை எடுத்துக்கொள்க:
$!A \land !B, \Gamma \Sequent \Delta$
என்ற இறுதித் தொடரணியுடன்
!!a{proof}~$\pi$ உள்ளது என்க.
$!A, !B, \Gamma \Sequent \Delta$
என்பதற்கான !!a{proof}~$\pi'$
உண்டு; மேலும் $\pheight{\pi'} \le
\pheight{\pi}$ என்பதைக் காட்ட வேண்டும்.
$n = \pheight{\pi}$ மீது
தொகுத்தறிந்து இதை நிறுவுவோம்.

$n = 0$ எனில், $\pi$ என்பது
$!A \land !B, \Gamma \Sequent \Delta$
என்ற அடிக்கோள் மட்டுமே.
\Log{G3c} இல் அடிக்கோள்கள்,
$!C$ அணுவாக இருக்கும்
$!C, \Gamma \Sequent \Delta, !C$
என்ற வடிவில் உள்ளன. ஆகவே
$!C$, $!A \land !B$ ஆக முடியாது.
$!A \land !B, \Gamma \Sequent \Delta$
ஓர் அடிக்கோள் எனில், ஏதோ அணு~$!C$,
$\Gamma$ விலும்~$\Delta$ விலும்
!!a{element} ஆக இருக்க வேண்டும்.
அப்போது $!A, !B, \Gamma \Sequent \Delta$
உம் ஓர் அடிக்கோள்; வேண்டிய
!!{proof}~$\pi'$ கிடைத்தது.
அதன் !!{height} உம்~$0$.

% SEGMENT OLP-0701-S04
$n>0$ எனில், அது ஓர் அனுமானத்தில்
முடிகிறது. உயரம்~$n$ க்கான
கூற்றை நிறுவ வேண்டும்; ஆனால்
சூழல் வேறுபட்டாலும் உயரம்~$<n$
உடைய எந்த !!{proof} க்கும்
அது மெய் என்று கருதலாம்.
வேறு சொற்களில்,
எந்த $k<n$, எந்த $\Pi$, $\Lambda$
ஆகியவற்றுக்கும்,
$!A \land !B, \Pi \Sequent \Lambda$
என்பதற்கு !!{height}~$k$
உடைய !!a{proof} இருந்தால்,
$!A, !B, \Pi \Sequent \Lambda$
என்பதற்கு !!{height}~$\le k$
உடைய !!a{proof} உண்டு.

இரு நிலைகள் உள்ளன:
இடப்புறம் உள்ள $!A \land !B$
முதன்மை வாய்பாடாகவோ,
அவ்வாறு இல்லாமலோ இருக்கும்.
அது முதன்மையானால்,
கடைசி அனுமானம்
\LeftR{\land}; உடனடித்
துணை-!!{proof}~$\pi_1$
$!A, !B, \Gamma \Sequent \Delta$
இல் முடிகிறது.
$\pheight{\pi_1} < \pheight{\pi}$
என்பதால் இந்நிலையில்
கூற்று உடனே கிடைக்கிறது.

% SEGMENT OLP-0701-S05
$!A \land !B$ முதன்மை
வாய்பாடல்ல எனில்,
வேறு ஓர் !!{formula}
முதன்மையாக இருக்கும்;
$!A \land !B$ கடைசி
அனுமானத்தின் சூழலில் இருக்கும்.
அப்போது கடைசி விதி~$R$
இன் உடனடித் துணை-!!{proof}s
மீது தொகுத்தறிதல் கருதுகோளைப்
பயன்படுத்தலாம். இதனால்
!!a{proof}~$\pi_1'$
அல்லது !!{proof}s~$\pi_1'$,
$\pi_2'$ கிடைக்கும்; அவற்றின்
!!{height} கள் $\pi$ இன்
உடனடித் துணை-!!{proof}s இன்
உயரங்களை மீறாது. $\pi_1'$,
$\pi_2'$ ஆகியவற்றின்
இறுதித் தொடரணிகளின்
இடச் சூழலில் $!A \land !B$
க்குப் பதிலாக $!A, !B$
இருக்கும். விதி~$R$ ஐ
மீண்டும் பயன்படுத்தி
நிறுவல்~$\pi'$ ஐப் பெறலாம்.
எடுத்துக்காட்டாக, $\pi$,
\LeftR{\lif} இல்
முடிகிறது என்க:
\[
\AxiomC{}
\RightLabel{$\pi_1$}
\Deduce$!A \land !B, \Gamma' \fCenter \Delta, !C$
\AxiomC{}
\RightLabel{$\pi_2$}
\Deduce$!A \land !B, \Gamma', !D \fCenter \Delta$
\RightLabel{\LeftR{\lif}}
\BinaryInf$!A \land !B, \Gamma', !C \lif !D \fCenter \Delta$
\DisplayProof
\]

% SEGMENT OLP-0701-S06
இங்கு இடப்புறம் உள்ள
$!C \lif !D$ முதன்மை
வாய்பாடு; இடச் சூழல்
$!A \land !B, \Gamma'$
(அதாவது $\Gamma = \Gamma', !C \lif !D$).
தொகுத்தறிதல் கருதுகோளால்,
முறையே $!A, !B, \Gamma' \Sequent
\Delta, !C$ மற்றும்
$!A, !B, \Gamma', !D \Sequent \Delta$
என்பவற்றுக்கான !!{proof}s
$\pi_1'$, $\pi_2'$ கிடைக்கும்.
இந்த இரு தொடரணிகளையும்
\LeftR{\lif} இன்
முன்கூற்றுகளாக மீண்டும்
பயன்படுத்தி~$\pi'$ ஐப்
பெறலாம்:
\[
\AxiomC{}
\RightLabel{$\pi_1'$}
\Deduce$!A, !B, \Gamma' \fCenter \Delta, !C$
\AxiomC{}
\RightLabel{$\pi_2'$}
\Deduce$!A, !B, \Gamma', !D \fCenter \Delta$
\RightLabel{\LeftR{\lif}}
\BinaryInf$!A, !B, \Gamma', !C \lif !D \fCenter \Delta$
\DisplayProof
\]
வரையறைப்படி,
\begin{align*}
\pheight{\pi} & = \max(\pheight{\pi_1}, \pheight{\pi_2}) + 1\\
\pheight{\pi'} & = \max(\pheight{\pi_1'}, \pheight{\pi_2'}) + 1
\intertext{மேலும் தொகுத்தறிதல் கருதுகோளின்படி,}
\pheight{\pi_1'} & \le \pheight{\pi_1}\\
\pheight{\pi_2'} & \le \pheight{\pi_2}
\intertext{எனவே,}
\pheight{\pi'} & \le \pheight{\pi}.
\end{align*}
\end{proof}

% SEGMENT OLP-0701-S07
\begin{prob}
\RightR{\lif} க்கு
\olref{lem:G3c-invert} இன்
நிறுவலைத் தருக.
\end{prob}

% SEGMENT OLP-0701-S08
\begin{lem}\ollabel{lem:invert-quant}
\RightR{\lforall},
\LeftR{\lexists}
ஆகிய விதிகள்,
பின்வரும் பொருளில்
!!{height}-பேணும்
மீளாக்கத்தக்கவை:
\begin{enumerate}
  \item $\Log{G3c} \Proves[n] \Gamma \Sequent \Delta,
  \lforall[x][!A(x)]$ எனில்,
  $\Log{G3c} \Proves[n] \Gamma \Sequent
  \Delta, !A(c)$; மேலும்
  \item $\Log{G3c} \Proves[n] \lexists[x][!A(x)], \Gamma \Sequent
  \Delta$ எனில்,
  $\Log{G3c} \Proves[n] !A(c), \Gamma \Sequent \Delta$.
\end{enumerate}
இங்கு முறையே $\Gamma$, $\Delta$,
$\lforall[x][!A(x)]$ அல்லது
$\lexists[x][!A(x)]$ ஆகியவற்றில்
வராத எந்த $c$ ஐயும்
எடுத்துக்கொள்ளலாம்.
\end{lem}

% SEGMENT OLP-0701-S09
\begin{proof}
  (1) ஐ நிறுவுவோம்;
  (2) பயிற்சியாக விடப்படுகிறது.
  வாதம் \olref{lem:G3c-invert}
  இன் நிறுவலைப் போன்றது.
  தொகுத்தறிதல் படியிலும்
  இரண்டு நிலைகள் உள்ளன.
  $\lforall[x][!A(x)]$
  முதன்மை வாய்பாடாக
  இருக்கும் நிலை எளிது:
  கடைசி $\RightR{\lforall}$
  அனுமானத்தின் முன்கூற்று
  வேண்டிய $\Gamma
  \Sequent \Delta, !A(a)$
  என்ற வடிவில் உள்ளது;
  அதில் முடியும் உடனடித்
  துணை-!!{proof}~$\pi_1$
  க்கு $\pheight{\pi_1}
  < \pheight{\pi}$.
  தனிமாறி நிபந்தனையால்
  $a$, $\Gamma$, $\Delta$
  அல்லது $\lforall[x][!A(x)]$
  இல் வராது.
  $\pi_1$ ஒழுங்கானது
  எனக் கொள்ளலாம்;
  ஆகவே
  \olref[qua]{lem:subst-regular}
  படி $\Subst{\pi_1}{c}{a}$
  என்பது அதே
  !!{height} உடைய
  $\Gamma \Sequent \Delta,
  !A(c)$ என்பதற்கான
  !!a{proof}.

% SEGMENT OLP-0701-S10
  இரண்டாம் நிலையில்,
  $\lforall[x][!A(x)]$
  கடைசி அனுமானத்தில்
  முதன்மையல்ல; வேறு
  வாய்பாடு முதன்மை.
  மீண்டும் ஓர் எடுத்துக்காட்டைப்
  பார்ப்போம். கடைசி அனுமானம்
  ~\RightR{\land} என்க.
  $\Delta$ வில் $!B \land !C$
  என்ற வடிவிலான ஓர்
  !!{formula} முதன்மையாக
  இருக்கும். !!{proof}~$\pi$
  இவ்வாறு முடிகிறது:
  \[
\AxiomC{}
\RightLabel{$\pi_1$}
\Deduce$\Gamma \fCenter \Delta', \lforall[x][!A(x)], !B$
\AxiomC{}
\RightLabel{$\pi_2$}
\Deduce$\Gamma \fCenter \Delta', \lforall[x][!A(x)], !C$
\RightLabel{\RightR{\land}}
\BinaryInf$\Gamma \fCenter \Delta', \lforall[x][!A(x)], !B \land !C$
\DisplayProof
\]
  இங்கு $\Delta = \Delta', !B \land !C$.
  $c$ என்பது $\Delta$ வில்
  வராத !!a{constant} என்க;
  அப்போது அது
  $\Delta', !B$
  இலோ $\Delta', !C$
  இலோ வராது.
  எனவே தொகுத்தறிதல்
  கருதுகோள் $\pi_1$,
  ~$\pi_2$ ஆகியவற்றுக்குப்
  பொருந்தும். இதனால்
  $\pheight{\pi_1'} \le
  \pheight{\pi_1}$
  மற்றும்
  $\pheight{\pi_2'} \le
  \pheight{\pi_2}$
  உடைய !!{proof}s~$\pi_1'$,
  $\pi_2'$ கிடைக்கும்.
  வேண்டிய $\Gamma \Sequent
  \Delta, !A(c)$
  என்பதன் !!{proof}~$\pi'$
  இதோ:
\[
\AxiomC{}
\RightLabel{$\pi_1'$}
\Deduce$\Gamma \fCenter \Delta', !A(c), !B$
\AxiomC{}
\RightLabel{$\pi_2'$}
\Deduce$\Gamma \fCenter \Delta', !A(c), !C$
\RightLabel{\RightR{\land}}
\BinaryInf$\Gamma \fCenter \Delta', !A(c), !B \land !C$
\DisplayProof
\]
  இதன் $\pheight{\pi'} =
  \max(\pheight{\pi_1'},
  \pheight{\pi_2'})+1
  \le \pheight{\pi}$.
\end{proof}

% SEGMENT OLP-0701-S11
வெட்டுநீக்கத் தேற்றத்தின் நிரூபணத்தில்
\Olref{lem:G3c-invert} முக்கியப் பங்கு வகிக்கிறது.
அதைக் கொண்டு பின்வருவதையும் காட்டலாம்:

\begin{prop}\ollabel{prop:G3c-cont-adm}
\Log{G1c} இன் சுருக்க விதிகளான
\LeftR{\Contraction}, \RightR{\Contraction}
ஆகியவை \Log{G3c} இல்
!!{height}-பேணும் வகையில் ஏற்கத்தக்கவை.
\end{prop}

% SEGMENT OLP-0701-S12
\begin{proof}
\RightR{\Contraction}, \LeftR{\Contraction}
ஆகிய இரண்டிற்கும் ஒரே நேரத்தில் வாதிடுவோம்.
$\pi$ என்பது \Log{G3c} இல்
$\Gamma \Sequent \Delta, !A, !A$
அல்லது $!A, !A, \Gamma \Sequent \Delta$
என்பதற்கான !!a{proof} என்க.
முறையே $\Gamma \Sequent \Delta, !A$
அல்லது $!A, \Gamma \Sequent \Delta$
என்பதற்கும் $\pheight{\pi'} \le
\pheight{\pi}$ உடைய \Log{G3c}-!!{proof}~$\pi'$
உண்டு என்பதைக் காட்ட வேண்டும்.
$n = \pheight{\pi}$ மீது தொகுத்தறிதல் செய்வோம்.

% SEGMENT OLP-0701-S13
$n=0$ எனில் $\pi$ ஓர் அடிக்கோளே.
$\Gamma \Sequent \Delta, !A, !A$
\Log{G3c} இன் அடிக்கோள் எனில்,
$\Gamma \Sequent \Delta, !A$ உம் அடிக்கோள்தான்:
ஏனெனில் அணு வாய்பாடான $!C$ ஒன்று
$\Gamma$, $\Delta$ இரண்டிலும்
!!a{element} ஆக இருக்கும்;
அல்லது $!A$ அணு வாய்பாடாகவும்
$\Gamma$ வின் !!a{element} ஆகவும் இருக்கும்.
முடிவுத் தொடர்கோவை $!A, !A, \Gamma \Sequent
\Delta$ எனிலும் இதே வாதம் பொருந்தும்.

% SEGMENT OLP-0701-S14
$n>0$ எனில் $\pi$ ஏதோ விதி~$R$ இல் முடிகிறது.
இந்தக் கடைசி அனுமானத்தில் $!A$ முதன்மையல்ல
எனில், \olref{lem:G3c-invert} இன் நிரூபணத்தைப் போல,
$\pi$ இன் உடனடி உட்-!!{proof}s மீது
தொகுத்தறிதல் கருதுகோளைப் பயன்படுத்தி,
கிடைத்த !!{proof}(s) மீது அதே விதி~$R$ ஐப்
பயன்படுத்தி !!a{proof}~$\pi'$ பெறலாம்.
இங்கு தொகுத்தறிதல் கருதுகோள் இதுதான்:
$k<n$ !!{height} உடைய
$\Pi \Sequent \Lambda, !D, !D$
அல்லது $!D, !D, \Pi \Sequent \Lambda$
என்பதற்கு !!a{proof} இருந்தால்,
முறையே $\Pi \Sequent \Lambda, !D$
அல்லது $!D, \Pi \Sequent \Lambda$
என்பதற்கு $\le k$ !!{height} உடைய
!!a{proof} உண்டு. (சூழல்கள் அல்லது
சுருக்கப்படும் !!{formula} ஆகியவை
$\Gamma$, $\Delta$, $!A$ விலிருந்து
வேறுபட்டிருந்தாலும் தொகுத்தறிதல் கருதுகோள்
பொருந்தும் என்பதைக் கவனிக்கவும்!)

எடுத்துக்காட்டாக $R$ என்பது $\RightR{\land}$
என்றும் $\pi$ இவ்வாறு முடிகிறது என்றும் கொள்வோம்:
\[
\AxiomC{}
\RightLabel{$\pi_1$}
\Deduce$\Gamma \fCenter \Delta', !B, !A, !A$
\AxiomC{}
\RightLabel{$\pi_2$}
\Deduce$\Gamma \fCenter \Delta', !C, !A, !A$
\RightLabel{\RightR{\land}}
\BinaryInf$\Gamma \fCenter \Delta', !B \land !C, !A, !A$
\DisplayProof
\]
இங்கு $\Delta = \Delta', !B \land !C$.
தொகுத்தறிதல் கருதுகோளால்,
முறையே $\Gamma \fCenter \Delta', !B, !A$
மற்றும் $\Gamma \fCenter \Delta', !C, !A$
என்பவற்றுக்கான $\pi_1'$, $\pi_2'$ என்னும்
!!{proof}s கிடைக்கின்றன; மேலும்
$\pheight{\pi_1'} \le  \pheight{\pi_1}$,
$\pheight{\pi_2'} \le \pheight{\pi_2}$.
வேண்டிய !!{proof}~$\pi'$ இதோ:
\[
\AxiomC{}
\RightLabel{$\pi_1'$}
\Deduce$\Gamma \fCenter \Delta', !B, !A$
\AxiomC{}
\RightLabel{$\pi_2'$}
\Deduce$\Gamma \fCenter \Delta', !C, !A$
\RightLabel{\RightR{\land}}
\BinaryInf$\Gamma \fCenter \Delta', !B \land !C, !A$
\DisplayProof
\]

% SEGMENT OLP-0701-S15
கடைசி அனுமானத்தில் $!A$ முதன்மை எனில்,
$\pi$ இன் உடனடி உட்-!!{proof}s மீது
மீளாக்கத் துணைமுடிவை
(\olref{lem:G3c-invert}) பயன்படுத்துவோம்;
பின்னர் தொகுத்தறிதல் கருதுகோளையும்,
மீண்டும் விதி~$R$ ஐயும் பயன்படுத்துவோம்.
எடுத்துக்காட்டாக, $\pi$ என்பது
$\Gamma \Sequent \Delta, !A, !A$ ஐ
நிறுவுகிறது என்றும்,
$!A \ident (!B \land !C)$ கடைசி
அனுமானத்தில் முதன்மை என்றும் கொள்வோம்.
அப்போது $\pi$ இவ்வாறு இருக்கும்:
\[
\AxiomC{}
\RightLabel{$\pi_1$}
\Deduce$\Gamma \fCenter \Delta, !B \land !C, !B$
\AxiomC{}
\RightLabel{$\pi_2$}
\Deduce$\Gamma \fCenter \Delta, !B \land !C, !C$
\RightLabel{\RightR{\land}}
\BinaryInf$\Gamma \fCenter \Delta, !B \land !C, !B \land !C$
\DisplayProof
\]
$\pi_1$ மீது \olref{lem:G3c-invert} ஐப்
பயன்படுத்தி, $\Gamma \Sequent \Delta, !B, !B$
என்பதற்கான !!a{proof}~$\pi_1'$ ஐப் பெறுகிறோம்.
(அதனால் $\Gamma \Sequent \Delta, !C, !B$
என்பதற்கான !!a{proof} உம் கிடைக்கும்;
ஆனால் அது இங்கு தேவையில்லை.)
அதேபோல் $\pi_2$ மீது
\olref{lem:G3c-invert} ஐப் பயன்படுத்தி,
$\Gamma \Sequent \Delta, !C, !C$
என்பதற்கான !!a{proof}~$\pi_2'$ ஐப் பெறுகிறோம்.
\RightR{\land} இன் மீளாக்கம்
!!{height}-பேணுவது என்பதால்,
$\pheight{\pi_1'} \le  \pheight{\pi_1} < \pheight{\pi}$
மற்றும்
$\pheight{\pi_2'} \le  \pheight{\pi_2} < \pheight{\pi}$.
எனவே $\pi_1'$, $\pi_2'$ மீது தொகுத்தறிதல்
கருதுகோள் பொருந்தும்: முறையே
$\Gamma \Sequent \Delta, !B$,
$\Gamma \Sequent \Delta, !C$
என்பவற்றுக்கான !!{proof}s~$\pi_1''$,
$\pi_2''$ உள்ளன; அவற்றுக்கு
$\pheight{\pi_1''} \le \pheight{\pi_1'}
\le \pheight{\pi_1}$ மற்றும்
$\pheight{\pi_2''} \le \pheight{\pi_2'}
\le \pheight{\pi_2}$.
ஆகவே !!{proof}~$\pi'$ இதுவாகும்:
\[
\AxiomC{}
\RightLabel{$\pi_1''$}
\Deduce$\Gamma \fCenter \Delta, !B$
\AxiomC{}
\RightLabel{$\pi_2''$}
\Deduce$\Gamma \fCenter \Delta, !C$
\RightLabel{\RightR{\land}}
\BinaryInf$\Gamma \fCenter \Delta, !B \land !C$
\DisplayProof
\]
இங்கு $\pheight{\pi'} \le \pheight{\pi}$.

% SEGMENT OLP-0701-S16
முதன்மை !!{formula} அளவையீட்டைக் கொண்டிருந்தால்,
அதற்குரிய விதிகளில் தனிக் கவனம் தேவை.

$!A \ident \lforall[x][!B(x)]$
வலப்புறத்தில் முதன்மை என்க.
அப்போது $\pi$ இவ்வாறு முடிகிறது:
\[
\AxiomC{}
\RightLabel{$\pi_1$}
\Deduce$\Gamma \fCenter \Delta, \lforall[x][!B(x)], !B(a)$
\RightLabel{\RightR{\lforall}}
\UnaryInf$\Gamma \fCenter \Delta, \lforall[x][!B(x)], \lforall[x][!B(x)]$
\DisplayProof
\]
\olref{lem:invert-quant} படி,
$\Gamma \fCenter \Delta,
\lforall[x][!B(x)], !B(a)$ இல்
வராத எந்த $c$ க்கும்,
$\Gamma \fCenter \Delta, !B(c), !B(a)$
என்பதற்கான $\le \pheight{\pi_1}$
!!{height} உடைய !!a{proof}~$\pi_1'$ உள்ளது.
$\pi_1'$ ஒழுங்கானது எனக் கொள்ளலாம்;
எனவே \olref[qua]{lem:subst-regular} படி
!!{proof} $\Subst{\pi_1'}{a}{c}$ என்பது
$\Gamma \fCenter \Delta, !B(a), !B(a)$
என்பதற்கான, $\le \pheight{\pi_1}$
!!{height} உடைய !!a{proof}.
இப்போது தொகுத்தறிதல் கருதுகோளால்
$\Gamma \fCenter \Delta, !B(a)$
என்பதற்கான !!a{proof}~$\pi_1''$ கிடைக்கிறது.
அதன் மீது \RightR{\lforall} விதியைப்
பயன்படுத்தி $\pi'$ ஐப் பெறுகிறோம்:
\[
\AxiomC{}
\RightLabel{$\pi_1''$}
\Deduce$\Gamma \fCenter \Delta, !B(a)$
\RightLabel{\RightR{\lforall}}
\UnaryInf$\Gamma \fCenter \Delta, \lforall[x][!B(x)]$
\DisplayProof
\]
இங்கு $\pheight{\pi'} \le \pheight{\pi}$.

% SEGMENT OLP-0701-S17
இப்போது $!A \ident \lexists[x][!B(x)]$
வலப்புறத்தில் முதன்மை என்க.
அப்போது $\pi$ இவ்வாறு உள்ளது:
\[
\AxiomC{}
\RightLabel{$\pi_1$}
\Deduce$\Gamma \fCenter \Delta, \lexists[x][!B(x)], \lexists[x][!B(x)], !B(t)$
\RightLabel{\RightR{\lexists}}
\UnaryInf$\Gamma \fCenter \Delta, \lexists[x][!B(x)], \lexists[x][!B(x)]$
\DisplayProof
\]
$\pi_1$ மீது தொகுத்தறிதல் கருதுகோள்
பொருந்தும். அதனால் $\Gamma \Sequent \Delta,
\lexists[x][!B(x)], !B(t)$ என்பதற்கான
!!a{proof}~$\pi_1'$ கிடைக்கும்.
(இங்கு மீளாக்கத் துணைமுடிவு தேவையில்லை
என்பதைக் கவனிக்கவும்!)
அப்போது !!{proof}~$\pi'$ இதுவாகும்:
\[
\AxiomC{}
\RightLabel{$\pi_1'$}
\Deduce$\Gamma \fCenter \Delta,  \lexists[x][!B(x)], !B(t)$
\RightLabel{\RightR{\lexists}}
\UnaryInf$\Gamma \fCenter \Delta, \lexists[x][!B(x)]$
\DisplayProof
\]
இங்கு $\pheight{\pi'} \le \pheight{\pi}$.
\end{proof}

% SEGMENT OLP-0701-S18
\begin{prob}
\LeftR{\Contraction} க்கான
\olref{prop:G3c-cont-adm} இன் நிரூபணத்தை
வரைந்து காட்டுக. $!A \ident (!B \land !C)$
மற்றும் $!A \ident (!B \lif !C)$
ஆகிய நிலைகளை விவரிக்கவும்.
\end{prob}

\begin{prob}
எஞ்சிய அளவையீட்டு நிலைகளில், அதாவது
$!A \ident \lforall[x][!B(x)]$
மற்றும் $!A \ident \lexists[x][!B(x)]$
எனும்போது, \LeftR{\Contraction} க்கான
\olref{prop:G3c-cont-adm} இன் நிரூபணத்தை
வரைந்து காட்டுக.
\end{prob}

\end{document}
